\section*{Capítulo \ref{ch:b3:exoplanets} --- Exoplanetas}

\begin{solution}{exo:b3:exoplanets:1}
(a)~$K = (2\pi G/P)^{1/3}m_p/\bigl((M_\odot + m_p)^{2/3}\sqrt{1 - e^2}\bigr)$.
Júpiter: $P = \qty{3.742e8}{s}$, $(2\pi G/P)^{1/3} = \num{1.039e-6}$ en
unidades SI, $K = \num{1.039e-6}\times\num{1.898e27}/(\num{1.990e30})^{2/3}/0.9988
= \qty{12.5}{m.s^{-1}}$. La Tierra: $P = \qty{3.155e7}{s}$, $K =
\qty{0.089}{m.s^{-1}}$, unos \qty{9}{cm.s^{-1}}. (b)~\qty{3.5}{km/h} $=
\qty{0.97}{m.s^{-1}}$ y \qty{0.35}{km/h} $= \qty{0.097}{m.s^{-1}}$.
La señal de Júpiter es trece veces la primera precisión; la de la Tierra es
inferior incluso a la segunda, y el ruido propio de la estrella se superpone.
\end{solution}

\begin{solution}{exo:b3:exoplanets:2}
(a)~$(71492/695700)^2 = \num{0.0106}$, una caída de \qty{1.06}{\percent};
$(6378/695700)^2 = \num{8.4e-5}$, 84 partes por millón. (b)~$R_\star/a$:
la Tierra $695700/\num{1.496e8} = \num{4.7e-3}$, aproximadamente una posibilidad entre 215;
Júpiter $695700/\num{7.785e8} = \num{8.9e-4}$, una posibilidad entre 1100.
\end{solution}

\begin{solution}{exo:b3:exoplanets:3}
$m_p\sin i \approx K M_\odot^{2/3}(P/2\pi G)^{1/3}$ con $P =
\qty{3.655e5}{s}$: $(P/2\pi G)^{1/3} = \num{9.55e4}$ en unidades SI, así que $m_p
\sin i = 59\times(\num{1.9884e30})^{2/3}\times\num{9.55e4} =
\qty{8.9e26}{kg}$, es decir, 0.47 masas de Júpiter o 149 masas terrestres (la
corrección de $m_p$ en $M_\star + m_p$ es inferior a una parte entre mil).
$a = \bigl(GM_\odot P^2/4\pi^2\bigr)^{1/3} = \qty{7.66e9}{m} =
\qty{0.051}{au}$.
\end{solution}

\begin{solution}{exo:b3:exoplanets:4}
(a)~$4GM/c^2 = 4\times\num{6.674e-11}\times\num{9.94e29}/(\num{2.998e8})^2 =
\qty{2953}{m}$; $(D_S - D_L)/(D_LD_S) = 1/D_S\times(D_S - D_L)/D_L =
1/(\qty{8}{kpc}) = 1/\qty{2.469e20}{m}$. $\theta_\mathrm E =
\sqrt{2953/\num{2.469e20}} = \qty{3.46e-9}{rad} = \qty{0.71}{mas}$. En la
lente, $R_\mathrm E = \theta_\mathrm E D_L = \num{3.46e-9}\times
\qty{1.234e20}{m} = \qty{4.27e11}{m} = \qty{2.85}{au}$. (b)~$A(1) =
3/\sqrt5 = 1.34$; $A(0.5) = 2.25/(0.5\times2.062) = 2.18$; $A(0.1) =
2.01/(0.1\times2.002) = 10.0$. (c)~$t_\mathrm E = \num{4.27e11}/\num{2e5}
= \qty{2.13e6}{s}$, unos 25 días: los eventos de microlente gravitacional duran semanas.
\end{solution}

\begin{solution}{exo:b3:exoplanets:5}
(a)~$a = (GM_\odot P^2/4\pi^2)^{1/3} = \qty{6.09e9}{m}$ (\qty{0.0407}{au}),
así que $a/R_\star = 8.75$. (b)~Con $P/\pi = \qty{22.9}{h}$, para $b = 0$:
\[
T_{14} = \frac P\pi\arcsin\frac{1.12}{8.75} = \qty{2.94}{h}, \qquad
T_{23} = \frac P\pi\arcsin\frac{0.88}{8.75} = \qty{2.31}{h}.
\]
Para $b = 0.7$: $\cos i = 0.7/8.75$, $i = \ang{85.4}$, $(a/R_\star)\sin i =
8.726$, y
\[
T_{14} = \frac P\pi\arcsin\frac{\sqrt{1.2544 - 0.49}}{8.726} = \qty{2.30}{h},
\qquad
T_{23} = \frac P\pi\arcsin\frac{\sqrt{0.7744 - 0.49}}{8.726} = \qty{1.40}{h}.
\]
(c)~La entrada y la salida ocupan $(T_{14} - T_{23})/T_{14} = \qty{21}{\percent}$
del tránsito cuando es central, \qty{39}{\percent} para $b = 0.7$; un
tránsito rasante, $b \to 1$, no tiene fondo plano y tiene forma de V.
\end{solution}

\begin{solution}{exo:b3:exoplanets:6}
$\bar\rho = M/\tfrac43\pi R^3$, con $\tfrac43\pi R_\oplus^3 =
\qty{1.087e21}{m^3}$: TRAPPIST-1~b $1.374\times\num{5.97e24}/(1.116^3\times
\num{1.087e21}) = \qty{5.43}{g.cm^{-3}}$; Kepler-10~b
\qty{5.60}{g.cm^{-3}}; WASP-39~b \qty{0.167}{g.cm^{-3}}; HD~209458~b
\qty{0.337}{g.cm^{-3}}. Los dos primeros son tan densos como la Tierra
(\qty{5.49}{g.cm^{-3}} con las mismas constantes) y son planetas rocosos; los dos últimos
son gigantes gaseosos, más ligeros que el agua, y WASP-39~b es ocho veces más ligero que
Júpiter.
\end{solution}

\begin{solution}{exo:b3:exoplanets:7}
Los extremos son $\gamma + K(1 + e\cos\omega)$ y $\gamma - K(1 -
e\cos\omega)$. Su diferencia es $2K = \qty{60.0}{m.s^{-1}}$, así que $K =
\qty{30.0}{m.s^{-1}}$; su media es $\gamma + Ke\cos\omega =
\qty{20.4}{m.s^{-1}}$, así que $Ke\cos\omega = \qty{8.0}{m.s^{-1}}$ y
$e\cos\omega = 0.27$. En una órbita excéntrica, la estrella se demora cerca
del apoastro y pasa deprisa por el periastro, así que la curva es asimétrica y el
punto medio de sus extremos no es su promedio temporal $\gamma$.
\end{solution}

\begin{solution}{exo:b3:exoplanets:8}
(a)~$g = GM/R^2 = \num{6.674e-11}\times\num{5.332e26}/(\num{9.144e7})^2 =
\qty{4.26}{m.s^{-2}}$; $H = kT/\bar m g = \num{1.381e-23}\times1100/(2.3
\times\num{1.661e-27}\times4.26) = \qty{934}{km}$. Con $R_\star =
\qty{6.533e8}{m}$, la profundidad es $(\num{9.144e7}/\num{6.533e8})^2 =
\num{0.0196}$. (b)~$2HR_p/R_\star^2 = 2\times\num{9.34e5}\times
\num{9.144e7}/(\num{6.533e8})^2 = \num{4.0e-4}$, 400 partes por millón;
dos alturas de escala, 800: un gigante gaseoso de gravedad baja es el blanco más fácil
para la \omterm{def:b3:exoplanets:transmission}{espectroscopía de transmisión}.
\end{solution}

\begin{solution}{exo:b3:exoplanets:9}
$p = 1.116\times6378/(0.1192\times695700) = 0.0858$, $p^2 = \num{7.37e-3}$.
$B_\lambda \propto 1/(e^{hc/\lambda kT} - 1)$ con $hc/\lambda k =
\qty{959}{K}$ a \qty{15}{\micro m}: la estrella $1/(e^{0.374} - 1) = 2.20$;
un lado diurno a \qty{400}{K}, $1/(e^{2.40} - 1) = 0.0998$, razón \num{0.0453},
profundidad \num{3.3e-4} (334 partes por millón); a \qty{500}{K}, $1/(e^{1.92} -
1) = 0.172$, razón \num{0.078}, profundidad \num{5.7e-4}. Cien kelvin
cambian la profundidad en un \qty{70}{\percent}: la medida de la profundidad fija la
temperatura del lado diurno con bastante precisión, y un lado diurno caliente significa que
se transporta poco calor al lado nocturno.
\end{solution}

\begin{solution}{exo:b3:exoplanets:10}
(a)~$S \propto L/d^2$, así que $d = \sqrt{0.005}\,\unit{au} = \qty{0.0707}{au}$;
$P = 2\pi\sqrt{a^3/GM} = \qty{1.33e6}{s} = \qty{15.4}{d}$.
(b)~$K = (2\pi G/P)^{1/3}M_\oplus/(0.2M_\odot)^{2/3} = \qty{0.75}{m.s^{-1}}$,
8.4 veces el efecto de la Tierra sobre el Sol: la estrella ligera se mueve más, y
el período corto también ayuda. (c)~$(6378/139140)^2 = \num{2.1e-3}$, 25 veces
el tránsito de la Tierra ante el Sol. Ambas señales son mayores; los tránsitos
se producen cada dos semanas, y la probabilidad de tránsito $R_\star/a =
\num{0.013}$ es tres veces la de la Tierra: las estrellas pequeñas ofrecen a los planetas rocosos en
sus zonas habitables la oportunidad de ser hallados y estudiados.
\end{solution}

\begin{solution}{exo:b3:exoplanets:11}
Con $r\,dr = -\mu\,d\mu$, la intensidad media es $\frac1\pi\int_0^1
I\,2\pi r\,dr = 2\int_0^1 I(\mu)\mu\,d\mu = 2\bigl[\tfrac12 - u_1(\tfrac12 -
\tfrac13) - u_2\int_0^1(1 - \mu)^2\mu\,d\mu\bigr]$. La última integral es
$\tfrac12 - \tfrac23 + \tfrac14 = \tfrac1{12}$, así que la media es $1 - u_1/3 -
u_2/6$, aquí $0.825$. Un planeta pequeño oculta $\pi p^2 I(\mu)$ de $\pi\times
0.825$: profundidad $p^2 I(\mu)/0.825$. En $r = 0$, $I = 1$: $1.212\,p^2$; en
$r = 0.5$ ($\mu = 0.866$), $I = 0.942$: $1.142\,p^2$; en $r = 0.9$ ($\mu =
0.436$), $I = 0.695$: $0.842\,p^2$. Un modelo de disco uniforme interpretaría
$1.212\,p^2$ como $p^2$ y sobrestimaría el radio en $\sqrt{1.212} = 1.10$,
\qty{10}{\percent}.
\end{solution}

\begin{solution}{exo:b3:exoplanets:12}
(a)~$R_\odot/a$: \num{9.3e-3} a \qty{0.5}{au}, \num{6.6e-3} a
\qty{0.7}{au}, \num{4.7e-3} a \qty{1}{au}. (b)~$42000\times0.11\times
\num{6.6e-3} = 31$ planetas transitan. (c)~Solo alrededor de un tercio de ellos
se encontraron: los demás dieron tránsitos demasiado poco profundos (84 partes por millón para
una Tierra) o demasiado escasos (uno o dos al año) para destacar sobre el ruido. Esa
fracción es la eficiencia de detección; sin corregirla, ni corregir por
la probabilidad de tránsito, los recuentos subestimarían cuán frecuentes son tales
planetas en un factor de unos 460.
\end{solution}

\begin{solution}{pb:b3:exoplanets:1}
\textbf{1.} $M_\star = 0.80\times\num{1.9884e30} = \qty{1.591e30}{kg}$, $P =
\qty{1.0368e6}{s}$: $a = (GM_\star P^2/4\pi^2)^{1/3} = \qty{1.4245e10}{m}
= \qty{0.0952}{au}$.

\textbf{2.} $v = 2\pi a/P = \qty{86.3}{km.s^{-1}}$.

\textbf{3.} $m_p\sin i = KM_\star^{2/3}(P/2\pi G)^{1/3} = 2.59\times\num{1.363e20}
\times\num{1.352e5} = \qty{4.77e25}{kg}$, \num{7.99} masas terrestres.

\textbf{4.} $a_\star = a\,m_p/M_\star = \num{1.4245e10}\times\num{4.77e25}/
\num{1.591e30} = \qty{427}{km}$; $2\pi a_\star/P = \qty{2.59}{m.s^{-1}}$,
la $K$ observada.

\textbf{5.} $2.59$ veces una precisión de \qty{1}{m.s^{-1}}, 26 veces una de
\qty{0.1}{m.s^{-1}}: el planeta está al alcance de los mejores instrumentos,
apenas al de la generación anterior.

\textbf{6.} $m_p\sin i \propto K\sqrt{1 - e^2}$: $\times0.980$, 7.83
masas terrestres, \qty{2}{\percent} menos.

\textbf{7.} $p = \sqrt{\num{8.38e-4}} = 0.02895$.

\textbf{8.} $R_\star = 0.76\times\qty{695700}{km} = \qty{528700}{km}$; $R_p =
\qty{15310}{km} = 2.40$ radios terrestres.

\textbf{9.} $a/R_\star = \num{1.4245e10}/\num{5.287e8} = 26.94$.

\textbf{10.} $T_{14} = (\qty{288}{h}/\pi)\arcsin(1.02895/26.94) = \qty{3.50}{h}$.

\textbf{11.} $\sqrt{(1 + p)^2 - b^2} = (a/R_\star)\sin(\pi T_{14}/P) = 26.94\times
\sin(\pi\times3.23/288) = 0.949$ (con $\sin i \approx 1$); $b^2 =
1.0587 - 0.9008 = 0.158$, $b = 0.40$.

\textbf{12.} $\cos i = b\,R_\star/a = 0.40/26.94 = 0.0148$, $i = \ang{89.15}$.

\textbf{13.} $R_\star/a = 1/26.94 = 0.037$: una probabilidad de cada 27.

\textbf{14.} $\sin i = 0.99989$: la masa verdadera es 7.99 masas terrestres, la \omterm{def:b3:exoplanets:rv}{masa mínima}
dentro de \num{1e-4}.

\textbf{15.} $\bar\rho = \num{4.77e25}/\bigl(\tfrac43\pi(\num{1.531e7})^3\bigr) =
\num{4.77e25}/\num{1.502e22} = \qty{3180}{kg.m^{-3}}$,
\qty{3.18}{g.cm^{-3}}.

\textbf{16.} El planeta es $2.40/1.74 = 1.38$ veces mayor que un planeta rocoso de
su masa, 2.6 veces su volumen. La roca por sí sola no puede formarlo: necesita una
envoltura de gas ligero, unos pocos por ciento de su masa en hidrógeno y helio,
o una fracción grande de agua.

\textbf{17.} Con 2.40 radios terrestres queda por encima del valle, en el pico de 2--3
radios terrestres: un \omterm{def:b3:exoplanets:classes}{subneptuno}.

\textbf{18.} $g = Gm_p/R_p^2 = \num{6.674e-11}\times\num{4.77e25}/(\num{1.531e7})^2 =
\qty{13.6}{m.s^{-2}}$.

\textbf{19.} $L = 4\pi R_\star^2\sigma T^4 = 4\pi(\num{5.287e8})^2\times\num{5.670e-8}
\times5000^4 = \qty{1.245e26}{W}$, 0.325 luminosidades solares.

\textbf{20.} $S/S_\oplus = 0.325/0.0952^2 = 35.9$: muy por dentro del borde interior de
la zona habitable.

\textbf{21.} $L/4\pi a^2 = \qty{4.88e4}{W.m^{-2}}$; $T_\mathrm{eq} =
\bigl(\num{4.88e4}\times0.7/(4\sigma)\bigr)^{1/4} = \qty{623}{K}$.

\textbf{22.} $H = kT/\bar m g = \num{1.381e-23}\times623/(2.3\times\num{1.661e-27}
\times13.6) = \qty{166}{km}$.

\textbf{23.} $2HR_p/R_\star^2 = 2\times\num{1.66e5}\times\num{1.531e7}/
(\num{5.287e8})^2 = \num{1.82e-5}$, 18 partes por millón; una banda situada a dos
alturas de escala, 36 partes por millón.

\textbf{24.} $H$ es $18/2.3 = 7.8$ veces menor, \qty{21}{km}: 2.3 partes por
millón por altura de escala, 4.6 para la banda. A diez partes por millón, la
atmósfera rica en hidrógeno muestra su banda de agua a más de tres veces
el ruido; una atmósfera de vapor aparece plana, como una nubosa o un planeta
sin aire.

\textbf{25.} \textbf{\qty[text-series-to-math=true]{3.2}{g/cm^3}}.
\end{solution}
